\section{Reinforcement Learning Training of Diffusion Models}
\label{sec:algs}

We now describe how RL algorithms can be used to train diffusion models. 
We present two classes of methods and show that each corresponds to a different mapping of the denoising process to the MDP framework.

\subsection{Problem Statement}
We assume a pre-existing diffusion model, which may be pretrained or randomly initialized. Assuming a fixed sampler, the diffusion model induces a sample distribution $p_\theta(\bx_0 \mid \bc)$.
The denoising diffusion RL objective is to maximize a reward signal $r$ defined on the samples and contexts:
\begin{equation*}
    \lrl(\theta) = \; \expect{
        \bc \sim p(\bc),~\bx_0 \sim p_\theta(\bx_0 \mid \bc)
    }{
        r(\bx_0, \bc)
    }
\end{equation*}
for some context distribution $p(\bc)$ of our choosing.


\subsection{Reward-Weighted Regression}
\label{sec:rwr}
To optimize $\lrl$ with minimal changes to standard diffusion model training, we can use the denoising loss $\ldiff$ (Equation~\ref{eq:denoising_objective}), but with training data $\bx_0 \sim p_\theta(\bx_0 \mid \bc)$ and an added weighting that depends on the reward $r(\bx_0, \bc)$.
\citet{lee2023aligning} describe a single-round version of this procedure for diffusion models, but in general this approach can be performed for multiple rounds of alternating sampling and training, leading to an online RL method.
We refer to this general class of algorithms as reward-weighted regression (RWR)~\citep{peters2007rwr}.

A standard weighting scheme uses exponentiated rewards to ensure nonnegativity,
\begin{align*}
    w_\methodrwr(\bx_0, \bc) = \frac{1}{Z} \exp\big(\beta r(\bx_0, \bc) \big),
\end{align*}
where $\beta$ is an inverse temperature and $Z$ is a normalization constant.
We also consider a simplified weighting scheme that uses binary weights,
\begin{align*}
    w_\sparse(\bx_0, \bc) = \mathds{1} \big[ r(\bx_0, \bc) \geq C \big],
\end{align*}
where $C$ is a reward threshold determining which samples are used for training.
In supervised learning terms, this is equivalent to repeated filtered finetuning on training data coming from the model.

Within the RL formalism, the RWR procedure corresponds to the following one-step MDP:
\begin{align*}
    \bs \triangleq \bc \hspace{2em}
    \ba \triangleq \bx_0 \hspace{2em}
    \pi(\ba \mid \bs) \triangleq p_\theta(\bx_0 \mid \bc) \hspace{2em}
    \rho_0(\bs) \triangleq p(\bc) \hspace{2em}
    R(\bs, \ba) \triangleq r(\bx_0, \bc)
\end{align*}
with a transition kernel $P$ that immediately leads to an absorbing termination state.
Therefore, maximizing $\lrl(\theta)$ is equivalent to maximizing the RL objective $\mathcal{J}_\text{RL}(\pi)$ in this MDP.

From RL literature, weighting a log-likelihood objective by $w_\methodrwr$ is known to approximately maximize $\mathcal{J}_{\text{RL}}(\pi)$ subject to a KL divergence constraint on $\pi$ \citep{nair2020accelerating}.
However, $\ldiff$ (Equation~\ref{eq:denoising_objective}) does not involve an exact log-likelihood --- it is instead derived as a variational bound on $\log p_\theta(\bx_0 \mid \bc)$.
Therefore, the RWR procedure applied to diffusion model training is not theoretically justified and only optimizes $\lrl$ very approximately.

\subsection{Denoising Diffusion Policy Optimization}
\label{sec:ddpo}

RWR relies on an approximate log-likelihood because it ignores the sequential nature of the denoising process, only using the final samples $\bx_0$.
In this section, we show how the denoising process can be reframed as a \emph{multi-step} MDP, allowing us to directly optimize $\lrl$ using policy gradient estimators.
This follows the derivation in \citet{fan2023sft}, who prove an equivalence between their method and a policy gradient algorithm where the reward is a GAN-like discriminator.
We present a general framework with an arbitrary reward function, motivated by our desire to optimize arbitrary downstream objectives (Section~\ref{sec:reward_functions}).
We refer to this class of algorithms as \methodfull{} (\methodours) and present two variants based on specific gradient estimators.

\textbf{Denoising as a multi-step MDP.}~
We map the iterative denoising procedure to the following MDP:
\begin{align*}
    \bs_t &\triangleq (\bc, t, \bx_t) &
    \pi(\ba_t \mid \bs_t) &\triangleq p_\theta(\bx_{t-1} \mid \bx_t, \bc) &
    P(\bs_{t+1} \mid \bs_t, \ba_t) &\triangleq \big( \delta_\bc, \delta_{t-1}, \delta_{\bx_{t-1}} \big) \\
    \ba_t &\triangleq \bx_{t-1} &
    \rho_0(\bs_0) &\triangleq \big(p(\bc), \delta_T, \mathcal{N}(\bzero, \bI)\big) &
    R(\bs_t, \ba_t) &\triangleq
        \begin{cases}
            r(\bx_0, \bc) & \text{if } t = 0 \\
            0 & \text{otherwise}
        \end{cases}
\end{align*}
in which $\delta_y$ is the Dirac delta distribution with nonzero density only at $y$.
Trajectories consist of $T$ timesteps, after which $P$ leads to a termination state.
The cumulative reward of each trajectory is equal to $r(\bx_0, \bc)$, so maximizing $\lrl(\theta)$ is equivalent to maximizing $\mathcal{J}_\text{RL}(\pi)$ in this MDP.


The benefit of this formulation is that if we use a standard sampler with $p_\theta(\bx_{t-1} \mid \bx_t, \bc)$ parameterized as in Equation~\ref{eq:sampler}, the policy $\pi$ becomes an isotropic Gaussian as opposed to the arbitrarily complicated distribution $p_\theta(\bx_0 \mid \bc)$ as it is in the RWR formulation.
This simplification allows for the evaluation of exact log-likelihoods and their gradients with respect to the diffusion model parameters.

\textbf{Policy gradient estimation.}~
With access to likelihoods and likelihood gradients, we can make direct Monte Carlo estimates of $\nabla_\theta \lrl$.
Like RWR, \methodours{} alternates collecting denoising trajectories $\{\bx_T, \bx_{T-1}, \dots, \bx_0\}$ via sampling and updating parameters via gradient descent.

The first variant of \methodours{}, which we call \methodreinforce{}, uses the score function policy gradient estimator, also known as the likelihood ratio method or REINFORCE \citep{williams1992simple,mohamed2020monte}:
\begin{align}
    \tag{\methodreinforce}
    \nabla_\theta \mathcal{J}_\text{DDRL} &= \expect{}{\; \sum_{t=0}^{T} \nabla_\theta \log p_\theta(\bx_{t-1} \mid \bx_t, \bc) \; r(\bx_0, \bc)} \label{eq:score_pg}
\end{align}
where the expectation is taken over denoising trajectories generated by the current parameters $\theta$.

However, this estimator only allows for one step of optimization per round of data collection, as the gradient must be computed using data generated by the current parameters.
To perform multiple steps of optimization, we may use an importance sampling estimator \citep{kakade2002approximately}:
\begin{align}
    \tag{\methodppo}
    \nabla_\theta \mathcal{J}_\text{DDRL} &= \expect{}{\; \sum_{t=0}^{T} \frac{p_\theta (\bx_{t-1} \mid \bx_t, \bc)}{p_{\theta_\text{old}} (\bx_{t-1} \mid \bx_t, \bc)} \; \nabla_\theta \log p_\theta(\bx_{t-1} \mid \bx_t, \bc) \; r(\bx_0, \bc)}
    \label{eq:is_pg}
\end{align}
where the expectation is taken over denoising trajectories generated by the parameters $\theta_\text{old}$.
This estimator becomes inaccurate if $p_\theta$ deviates too far from $p_{\theta_\text{old}}$, which can be addressed using trust regions \citep{schulman2015trpo} to constrain the size of the update.
In practice, we implement the trust region via clipping, as in proximal policy optimization \citep{schulman2017ppo}.